% Part: computability
% Chapter: recursive-functions
% Section: other-recursions

\documentclass[../../../include/open-logic-section]{subfiles}

\begin{document}

\olfileid[fr]{cmp}{rec}{ore}
\olsection{Autres formes de récurrence}

Le codage des couples et des suites permet de justifier des
formes moins ordinaires, mais utiles, de récurrence primitive.
Il est par exemple souvent commode de définir deux fonctions
à la fois, comme ici :
\begin{align*}
h_0(\vec x, 0) & = f_0(\vec x) \\
h_1(\vec x, 0) & = f_1(\vec x) \\
h_0(\vec x, y+1) & = g_0(\vec x, y, h_0(\vec x, y), h_1(\vec x, y)) \\
h_1(\vec x, y+1) & = g_1(\vec x, y, h_0(\vec x, y), h_1(\vec x, y))
\end{align*}
C'est une \emph{récurrence mutuelle}. On peut également
définir $h(\vec x,y+1)$ à partir de \emph{toutes} les valeurs
$h(\vec x,0),\dots,h(\vec x,y)$, par exemple :
\begin{align*}
h(\vec x, 0) & = f(\vec x) \\
h(\vec x, y+1) & = g(\vec x, y, \tuple{h(\vec x, 0), \dots, h(\vec x, y)}).
\end{align*}
Le schéma suivant exprime la même idée de manière plus concise :
\[
h(\vec x, y) = g(\vec x, y, \tuple{h(\vec x, 0), \dots, h(\vec x, y-1)})
\]
Lorsque $y=0$, le dernier argument de~$g$ est, par
convention, la suite vide. Dans les deux formulations, le
calcul de la valeur suivante peut utiliser toute la suite
des valeurs déjà obtenues. On parle de \emph{récurrence
sur la suite des valeurs}. Elle permet notamment de
justifier des définitions de la forme
\begin{align*}
h(\vec x, y) & = \begin{cases}
  g(\vec x, y, h(\vec x, k(\vec x, y))) & \text{si $k(\vec x, y) < y$} \\
  f(\vec x) & \text{sinon}
\end{cases}
\end{align*}
La valeur de $h$ en~$y$ peut donc dépendre de sa valeur
en \emph{n'importe quel} rang antérieur donné par~$k$.

\begin{prob}
  Définir par récurrence sur la suite des valeurs la fonction
  reste~$r(x,y)$. Pour $x,y$ naturels et $y>0$, $r(x,y)$
  est l'unique entier inférieur à~$y$ tel que
  $x=z\times y+r(x,y)$ pour un certain~$z$.
  On convient que $r(x,0)=0$.
\end{prob}

Réfléchissez à la manière de retrouver ces fonctions au
moyen de la récurrence primitive ordinaire. Une dernière
variante autorise aussi la modification des
\emph{paramètres} au cours du calcul :
\begin{align*}
h(\vec x, 0) & = f(\vec x) \\
h(\vec x, y+1) & = g(\vec x, y, h(k(\vec x), y))
\end{align*}
Ici $k(\vec x)$ fournit un nouveau tuple de paramètres de
même arité que~$\vec x$. On suppose $f$, $g$ et chaque composante de $k$ récursives primitives.
Cette variante se ramène elle aussi à la récurrence primitive ordinaire. La construction demande
un peu de soin : pour vous guider, déroulez le calcul à la main.

\end{document}
